\documentclass[11pt]{article}
\usepackage[margin=1in]{geometry}
\usepackage{hyperref}
\usepackage{microtype}
\usepackage{enumitem}
\usepackage{titlesec}
\usepackage{cite}
\setlength{\parskip}{0.5em}
\setlength{\parindent}{0pt}
\title{A Survey on retrieval augmented generation}
\author{Generated by Agentic AutoSurvey \\\\ (multi-agent reimplementation of arXiv:2509.18661)}
\date{May 18, 2026}
\begin{document}
\maketitle

\begin{abstract}
Retrieval-Augmented Generation (RAG) has become a cornerstone for enhancing Large Language Models (LLMs), grounding their outputs in external knowledge to mitigate hallucination and provide up-to-date information. This survey presents a comprehensive and structured synthesis of the RAG landscape. We organize our review into six critical sub-topics: (1) RAG System Innovations and Applications, detailing novel architectures and real-world use cases; (2) RAG for Medical and Healthcare LLMs, a high-stakes domain; (3) RAG Fundamentals and Evaluation, covering core mechanisms and metrics; (4) the crucial challenges of Security and Privacy of RAG; (5) the emerging paradigm of Graph Retrieval-Augmented Generation; and (6) a meta-analysis of existing Retrieval-Augmented Generation Surveys. Distinct from prior works, our survey adopts a holistic pipeline perspective, examining the intricate interplay between the retrieval, augmentation, and generation stages. By synthesizing across these domains, we identify key research tensions, highlight underexplored open problems, and chart a path toward more robust, trustworthy, and efficient RAG systems.
\end{abstract}

\section{Introduction}
Large Language Models (LLMs) have become foundational to modern artificial intelligence, demonstrating remarkable capabilities in natural language understanding and generation. Despite their power, LLMs suffer from inherent limitations stemming from their reliance on the static knowledge encapsulated within their parameters during training. This leads to critical issues such as factual inaccuracies, or "hallucinations"; an inability to access information beyond their training data cutoff; and a lack of transparency, making it difficult to trace the provenance of their generated outputs . These challenges severely restrict their reliability in knowledge-intensive, real-world applications where currency and factual grounding are non-negotiable.

Retrieval-Augmented Generation (RAG) has emerged as a leading paradigm to mitigate these fundamental shortcomings \cite{ref1}. The core principle of RAG is to synergistically combine the generative power of LLMs with the vast, dynamic knowledge stored in external databases. By first retrieving relevant documents or data snippets and then conditioning the generation process on this retrieved context, RAG systems can produce more accurate, timely, and attributable responses \cite{ref2}. This approach not only enhances factual accuracy but also offers a practical mechanism for updating an LLM's knowledge base without the prohibitive cost of full retraining. The field has rapidly evolved beyond the initial retrieve-and-generate framework \cite{ref1}with recent work exploring more sophisticated methods like active, multi-step retrieval \cite{ref3}and extreme context compression \cite{ref4}. This survey synthesizes this burgeoning literature, charting the evolution of RAG from its core principles to its diverse applications and the critical challenges that define the research frontier.

We structure our review to provide a comprehensive overview of the RAG landscape. We begin by covering **RAG Fundamentals and Evaluation**, which explores the core architectural components and the critical need for robust assessment frameworks to measure pipeline performance \cite{ref5}. Next, we survey **RAG System Innovations and Applications**, detailing advancements that enhance efficiency and tackle complex problems like multi-hop queries \cite{ref6}. We then delve into a specialized retrieval modality in **Graph Retrieval-Augmented Generation**, examining techniques that leverage structured graph data instead of unstructured text \cite{ref7}. Subsequently, we explore a high-stakes domain in **RAG for Medical and Healthcare LLMs**, highlighting efforts to improve diagnostic accuracy and safety in clinical settings . Finally, we address a crucial and often overlooked aspect in **Security and Privacy of RAG**, discussing vulnerabilities such as knowledge poisoning and data leakage that threaten the trustworthiness of these systems .

\section{RAG System Innovations and Applications}
As foundational Retrieval-Augmented Generation (RAG) systems mature, research is shifting to address their inherent architectural limitations and expand their applicability to more complex tasks. A key focus of recent work is moving beyond simple, independent document retrieval to tackle challenges such as inefficient context utilization, inadequate reasoning over multiple information sources, and the need for robust domain-specific adaptation. These efforts aim to enhance the efficiency, contextual awareness, and reasoning capabilities of RAG systems, pushing them toward more sophisticated real-world applications.

\subsection*{Enhancing Retrieval with Structured Knowledge}

A significant thrust in advancing RAG involves enriching the retrieval process with structured knowledge, moving away from the "flat" data representations that can lead to fragmented and contextually poor answers \cite{ref8}. One approach is to incorporate graph structures directly into the text indexing and retrieval pipeline. LightRAG, for instance, employs a dual-level retrieval system over a graph-indexed corpus, designed to capture complex inter-dependencies between documents and provide more comprehensive, synthesized information to the language model \cite{ref8}. This methodology is echoed in domain-specific applications, such as recommender systems, where standard RAG can suffer from the same LLM-inherent issues of hallucination and outdated knowledge \cite{ref9}. By leveraging domain-specific Knowledge Graphs (KGs) as the external knowledge source, RAG systems can be grounded in factual, up-to-date, and relevant structured information, thereby improving the quality and reliability of recommendations \cite{ref9}. Both approaches converge on the principle that modeling explicit relationships between knowledge snippets—whether as a general text graph or a formal KG—is crucial for deeper contextual understanding and more accurate generation.

\subsection*{Innovations in System Architecture and Evaluation}

Beyond improving the knowledge source, another line of work focuses on innovating the core architecture of the retriever-generator interface and rigorously evaluating its limitations. A primary bottleneck in RAG is the limited context window of LLMs, which restricts the amount of retrieved text that can be processed. To address this, the xRAG framework proposes a form of extreme context compression by re-purposing document embeddings \cite{ref4}. Instead of passing retrieved text, xRAG fuses the dense vector embedding of a document directly into the LLM's representation space via a lightweight, trainable "modality bridge." This effectively compresses the retrieved context into a single token-equivalent representation, freeing up the context window while keeping both the retriever and LLM frozen \cite{ref4}.

While architectural innovations aim to improve system capabilities, a parallel effort seeks to identify and benchmark critical failure modes. A notable weakness of many existing RAG systems is their inability to answer "multi-hop" queries, which require retrieving and reasoning over multiple, disparate pieces of evidence \cite{ref6}. The MultiHop-RAG benchmark was developed specifically to address this gap, providing a dataset and evaluation framework focused on such queries. This work demonstrates that current systems struggle with this complex reasoning task and provides a crucial tool for measuring future progress \cite{ref6}.

\subsection*{Open Problems and Future Directions}

Together, these works highlight several open challenges. While graph-based retrieval is powerful, the automatic construction, maintenance, and efficient querying of large-scale knowledge graphs remain significant engineering hurdles . Furthermore, radical compression techniques like xRAG introduce a trade-off between context efficiency and information fidelity; it is an open question how much crucial, fine-grained detail is lost when text is replaced entirely by its embedding \cite{ref4}. Finally, with the formalization of challenges like multi-hop reasoning, the field now requires novel architectures that can explicitly perform iterative retrieval and synthesis, moving beyond the single-shot retrieve-and-generate paradigm \cite{ref6}. The successful adaptation of RAG for specialized domains like recommender systems also points toward a future of highly customized RAG pipelines tailored to the unique data structures and reasoning patterns of specific fields \cite{ref9}.

\section{RAG for Medical and Healthcare LLMs}
\subsection*{RAG for Medical and Healthcare LLMs}

The application of Large Language Models (LLMs) in high-stakes domains like healthcare is fraught with challenges, primarily the risk of generating factually incorrect "hallucinations," reliance on outdated training data, and a general lack of transparency . Retrieval-Augmented Generation (RAG) has emerged as a critical technique to mitigate these risks by grounding LLM responses in external, authoritative knowledge sources. Research in this area focuses on adapting, evaluating, and enhancing RAG systems to meet the stringent accuracy and safety requirements of clinical applications.

A significant body of work aims to systematically understand and evaluate the current landscape of medical RAG. Systematic reviews and meta-analyses confirm that RAG consistently improves the performance of LLMs in biomedical applications compared to baseline models \cite{ref10}. However, these reviews also highlight a lack of standardized methodologies, datasets, and evaluation frameworks, which complicates cross-study comparisons and the establishment of best practices \cite{ref11}. To address this gap, dedicated benchmarks are being developed. For instance, the Medical Information Retrieval-Augmented Generation Evaluation (MIRAGE) benchmark was proposed to systematically assess the various components of a RAG pipeline across multiple medical question-answering datasets \cite{ref12}. This work underscores that there is no one-size-fits-all solution; the optimal RAG configuration often depends on the specific medical task \cite{ref12}.

Empirical studies provide strong evidence for RAG's utility in specific clinical scenarios. A large-scale assessment across ten different LLMs demonstrated that a RAG system integrating specialized knowledge from local and international medical guidelines could significantly improve accuracy, consistency, and safety in determining surgical fitness \cite{ref13}. The study found that a GPT-4-based RAG model achieved near-human performance when retrieving from high-quality international guidelines, reinforcing the principle that the quality of the retrieved knowledge is paramount \cite{ref13}.

Despite these successes, researchers recognize the limitations of standard, heuristic-based RAG. Such models may be inadequate for complex diagnostic tasks, particularly for diseases with similar manifestations that require nuanced reasoning beyond simple information retrieval \cite{ref14}. This has motivated the development of more sophisticated RAG architectures. One such approach, MedRAG, enhances the retrieval process by incorporating a knowledge graph (KG) to elicit a multi-step reasoning process. By constructing a hierarchical KG of diseases, manifestations, and treatments, MedRAG aims to improve diagnostic accuracy and specificity, especially when dealing with privacy-sensitive Electronic Health Records (EHRs) \cite{ref14}.

Looking forward, several open problems remain. The field requires further development of robust, standardized evaluation benchmarks to guide progress \cite{ref12}. There is a pressing need for clear clinical development guidelines to ensure the effective and safe deployment of these systems in practice \cite{ref10}. Furthermore, future work must move beyond simple retrieval to incorporate more complex reasoning mechanisms, such as those enabled by knowledge graphs, to tackle the full spectrum of clinical challenges \cite{ref14}. The generalizability of RAG models across different medical specialties and tasks also warrants deeper investigation \cite{ref13}.

\section{RAG Fundamentals and Evaluation}
Retrieval-Augmented Generation (RAG) has emerged as a foundational technique for mitigating the factual inaccuracies and knowledge limitations inherent in large language models (LLMs) . While LLMs store a vast amount of knowledge in their parameters, their ability to access and manipulate this information precisely is limited, often leading to hallucinations, especially in knowledge-intensive domains \cite{ref1}. RAG addresses this by coupling the parametric memory of an LLM with a non-parametric external knowledge source, such as a textual database, enabling the model to ground its responses in verifiable information . This not only improves factual accuracy but also provides provenance for generated claims and allows for the model's knowledge to be updated without costly retraining \cite{ref1}.

The seminal RAG architecture established a "retrieve-and-generate" pipeline, where a retriever first fetches a set of relevant documents from a knowledge corpus based on the user's input, and a generator then produces an answer conditioned on both the input and the retrieved context \cite{ref1}. This single-step retrieval approach proved highly effective, achieving state-of-the-art results on a range of knowledge-intensive NLP tasks \cite{ref1}. However, this static retrieval paradigm has been identified as a significant limitation, particularly for tasks requiring the generation of long, multi-faceted texts where the information needs may evolve throughout the generation process \cite{ref3}. To overcome this, more dynamic approaches have been proposed. Active retrieval augmented generation, for instance, reframes the problem by allowing the model to make sequential decisions about *when* and *what* to retrieve during generation, treating retrieval as an adaptive policy rather than a one-off preliminary step \cite{ref3}.

As RAG systems have grown more complex, their evaluation has become a critical challenge \cite{ref5}. The performance of a RAG pipeline is contingent on multiple factors, and a failure in the final output could stem from either the retrieval or the generation component. For example, the retriever might fail to identify relevant context, or the generator might fail to faithfully utilize the provided information \cite{ref5}. This multi-dimensional nature of RAG performance necessitates specialized evaluation frameworks. To this end, methodologies like RAGAs (Retrieval Augmented Generation Assessment) have been developed to offer a reference-free, automated evaluation of RAG pipelines \cite{ref5}. Such frameworks decompose the evaluation into component-level metrics, assessing the retriever on dimensions like context relevance and the generator on its faithfulness to the context and the relevance of its answer \cite{ref5}.

Despite these advancements, several open problems remain. For active retrieval methods, defining and learning optimal retrieval policies without extensive supervision is an ongoing research direction \cite{ref3}. Furthermore, while frameworks like RAGAs represent a significant step forward, the development of comprehensive, low-cost, and reliable evaluation metrics for the nuanced qualities of RAG systems remains an active area of investigation \cite{ref5}. Understanding and improving the intricate interplay between the retriever and generator, especially the generator's robustness to noisy or conflicting retrieved information, is another key frontier for future work.

\section{Security and Privacy of RAG}
While Retrieval-Augmented Generation (RAG) is designed to enhance the factuality and timeliness of Large Language Models (LLMs), its reliance on an external knowledge corpus introduces a significant new attack surface \cite{ref15}. Recent work has begun to systematically explore the security and privacy vulnerabilities inherent in this architecture, demonstrating that the retrieval database can be a potent vector for compromising the system's integrity, availability, and confidentiality. These studies collectively shift the security focus from the LLM alone to the tightly coupled retriever-generator system, revealing threats that are unique to the RAG paradigm.

\subsection*{Attacks on Corpus Integrity and Availability}

A primary category of threats involves an adversary manipulating the knowledge corpus to degrade the RAG system's performance. These data poisoning attacks assume the adversary can inject a small number of malicious documents into the retrieval database, a plausible scenario in systems that index user-generated content or the open web. The goals of these attacks, however, can differ significantly. One line of work focuses on "knowledge corruption," where poisoned documents are crafted to be retrieved for specific queries and contain subtle misinformation that the LLM then synthesizes into a plausible but incorrect final answer \cite{ref15}. This represents an integrity attack, aiming to erode user trust by making the RAG system a source of falsehoods.

In contrast, other research demonstrates attacks on system availability through a technique termed "jamming" \cite{ref16}. Here, an adversary inserts "blocker" documents designed not to misinform, but to prevent the RAG system from answering a targeted query altogether. These documents are engineered to be highly relevant to the query but contain content that triggers the LLM's safety filters or leads it to conclude that it lacks sufficient information, effectively causing a denial-of-service for specific topics \cite{ref16}. Both corruption and jamming attacks have been shown to be highly effective, often requiring only a single malicious document to compromise a target query, highlighting the fragility of RAG systems that operate on untrusted data sources .

\subsection*{Privacy Risks and Trade-offs}

Complementing the focus on malicious data injection, another critical dimension of RAG security is the privacy of the retrieval corpus itself, especially when it contains proprietary or sensitive information. Research demonstrates that RAG systems are vulnerable to novel privacy leakage attacks, where an adversary can craft queries to infer or directly extract sensitive contents from the private documents used for retrieval \cite{ref17}. This exposes a new privacy risk specific to the RAG architecture, as the generation process is explicitly conditioned on potentially private data.

Interestingly, the interaction between RAG and privacy is not entirely negative. The same study reveals a nuanced trade-off: while RAG introduces risks for the retrieval data, it can simultaneously mitigate the leakage of sensitive information memorized by the base LLM from its own pre-training data \cite{ref17}. By grounding the model's output in a specific, retrieved context, RAG can reduce the likelihood of the LLM divulging memorized personal information or copyrighted text. This duality presents a significant open problem: securing the RAG pipeline involves managing a complex balance between protecting the private retrieval corpus and preventing leakage from the foundational model. Developing robust defense mechanisms that can sanitize malicious documents, verify information sources, and navigate this privacy trade-off remains a critical and largely unexplored area for future research.

\section{Graph Retrieval-Augmented Generation}
While standard Retrieval-Augmented Generation (RAG) has proven effective for augmenting Large Language Models (LLMs) with external knowledge, it primarily operates on unstructured, "flat" text corpora \cite{ref18}. This approach struggles to capture and leverage the complex, inherent relationships between entities, a critical limitation when dealing with specialized domains where knowledge is often highly structured and interconnected. Graph RAG emerges as a powerful paradigm to address this gap by replacing flat document retrieval with retrieval from graph-structured knowledge bases, enabling LLMs to reason over complex relational data . This shift is motivated by the need to overcome key challenges in traditional RAG, including difficulties in understanding complex queries, integrating knowledge from distributed sources, and achieving system efficiency in professional contexts \cite{ref18}.

The core methodological innovation in Graph RAG is the integration of graph traversal and representation learning with the LLM's generative process. Instead of retrieving independent text chunks, these systems first identify and extract a relevant subgraph in response to a user's query \cite{ref7}. This retrieval step often employs a combination of graph neural networks (GNNs) and language models to semantically align the natural language query with the graph's nodes and edges. For instance, the G-Retriever framework is designed to support conversational question answering over textual graphs by first identifying the most salient subgraph and then presenting it as context to the LLM \cite{ref7}. This retrieved subgraph, which encapsulates the most relevant entities and their relationships, is then typically linearized or summarized into a textual format that can be injected into the LLM's prompt. This allows the model to generate responses that are not only factually grounded in the knowledge base but also structurally aware of the relationships between different pieces of information \cite{ref18}.

Empirically, Graph RAG frameworks have demonstrated significant promise for domain-specific customization, particularly for question-answering tasks over real-world textual graphs \cite{ref7}. A key advantage of this approach is enhanced interpretability; by highlighting the specific subgraph used to generate an answer, these systems can provide users with direct, visual evidence for the model's reasoning, a feature largely absent in conventional RAG \cite{ref7}. This makes Graph RAG particularly suitable for applications in science, finance, and enterprise knowledge management, where verifying the source and context of information is crucial \cite{ref18}.

Despite these advances, significant open problems remain. Scalability and efficiency are primary concerns, as retrieving optimal subgraphs from massive, real-world knowledge graphs in real-time is a formidable computational challenge \cite{ref18}. Another major hurdle is bridging the semantic gap between ambiguous natural language queries and the formal structure of the graph, especially for multi-hop questions that require complex relational reasoning . Future work may focus on developing more sophisticated graph-aware retrieval algorithms, exploring techniques for integrating knowledge from heterogeneous and distributed graph sources, and extending the Graph RAG paradigm beyond question answering to more complex generative tasks.

\section{Retrieval-Augmented Generation Surveys}
Recent survey literature converges on Retrieval-Augmented Generation (RAG) as a critical paradigm for addressing fundamental limitations inherent in Large Language Models (LLMs) . The core problem, as identified across these surveys, is the static nature of LLM training data, which renders them prone to several well-documented failures. These include the generation of factually incorrect information, or "hallucination" \cite{ref2}; an inability to access information beyond their training cutoff, leading to outdated responses ; and an opaque, untraceable reasoning process that undermines trust and verifiability \cite{ref2}. RAG directly confronts these issues by synergistically merging the parametric knowledge of an LLM with non-parametric knowledge retrieved from external, dynamic data sources \cite{ref2}. This process of grounding generation in external evidence is positioned as a key strategy for enhancing the accuracy, credibility, and robustness of AI-generated content \cite{ref19}while also providing a mechanism for continuous knowledge updates without costly retraining \cite{ref2}.

The surveys conceptualize RAG's architecture as an evolving pipeline, moving from simple, static workflows to more dynamic, intelligent systems. The foundational RAG model consists of a retriever that fetches relevant information from a knowledge base in response to a query, and a generator (the LLM) that synthesizes an answer based on both the original query and the retrieved context. This basic structure is a common touchstone in comprehensive reviews of the field . However, a key point of evolution and implicit disagreement in the literature concerns the sufficiency of this model. More recent analysis argues that such "traditional RAG" systems are constrained by their rigid, pre-defined workflows and lack the adaptability required for complex, multi-step reasoning tasks \cite{ref20}. This critique gives rise to the concept of "Agentic RAG," which is framed as the next frontier. In this advanced paradigm, an autonomous agent dynamically plans and executes retrieval and generation steps, potentially refining queries, reasoning over retrieved documents, and adapting its strategy to the complexity of the task at hand \cite{ref20}.

Despite this architectural evolution, there is a strong consensus on the empirical benefits of the RAG approach. The literature consistently finds that by grounding outputs in external, verifiable data, RAG significantly improves factual accuracy and mitigates hallucinations in knowledge-intensive applications . It is broadly recognized as an effective method for integrating real-time data and domain-specific knowledge, thereby keeping models current and specialized .

Nonetheless, significant open problems remain. The primary limitation of simpler RAG models is their fragility; the quality of the final output is critically dependent on the success of the initial retrieval step. If the retriever fails, the entire system can fail. The push towards Agentic RAG is a direct response to this limitation, seeking to build more resilient and adaptive systems \cite{ref20}. Future work highlighted by these surveys points towards developing more sophisticated agentic controllers, improving the synergy between retrieval and generation components, and enhancing the system's ability to handle noisy, conflicting, or long-tail data from diverse external sources .
\section{Conclusion and Future Directions}
Retrieval-Augmented Generation has rapidly evolved from a foundational technique for mitigating LLM hallucinations \cite{ref1}into a diverse and critical subfield of AI . Our synthesis reveals a clear trajectory from the initial static retrieve-and-generate paradigm to more dynamic and specialized architectures. Innovations are pushing the boundaries of RAG's capabilities, enabling active, multi-step information gathering \cite{ref3}and reasoning over complex, multi-hop queries \cite{ref6}. Simultaneously, research is exploring novel data structures, moving beyond flat text to leverage the rich relational information in knowledge graphs, thereby enhancing contextual understanding and domain-specific customization .

This increasing sophistication is mirrored in RAG's application to high-stakes domains like healthcare, where it is being systematically evaluated to ensure the delivery of safe, accurate, and up-to-date medical information . However, as RAG systems become more powerful and integrated, they introduce new and critical challenges. The external knowledge base, a core component of RAG, creates novel attack surfaces for knowledge poisoning and data manipulation \cite{ref15}. Furthermore, the interaction between the retriever and generator can expose sensitive information from private databases, creating significant privacy risks that are still not fully understood \cite{ref17}. Addressing these vulnerabilities requires a parallel advancement in evaluation methodologies, moving beyond simple accuracy metrics to create robust frameworks that can assess the multiple dimensions of a RAG pipeline's performance and resilience \cite{ref5}.

Looking forward, the continued progress of RAG hinges on addressing several key areas. Future research will develop unified benchmarks that jointly evaluate retrieval quality, generation faithfulness, and security under adversarial conditions. The next frontier of RAG systems will integrate heterogeneous knowledge sources, combining unstructured text, structured graphs, and other modalities to support more complex reasoning. The paradigm will shift further towards autonomous, agentic architectures that can proactively seek, verify, and synthesize information over extended interactions \cite{ref20}. Finally, developing novel context compression and representation techniques is crucial for scaling RAG to massive knowledge bases while maintaining low-latency performance \cite{ref4}.

\begin{thebibliography}{99}
\bibitem{ref1} Patrick Lewis, Ethan Perez, Aleksandara Piktus, F. Petroni, Vladimir Karpukhin, Naman Goyal, et al.. Retrieval-Augmented Generation for Knowledge-Intensive NLP Tasks. \emph{Neural Information Processing Systems}, 2020. \url{https://www.semanticscholar.org/paper/659bf9ce7175e1ec266ff54359e2bd76e0b7ff31}
\bibitem{ref2} Yunfan Gao, Yun Xiong, Xinyu Gao, Kangxiang Jia, Jin Pan, Yuxi Bi, et al.. Retrieval-Augmented Generation for Large Language Models: A Survey. \emph{arXiv.org}, 2023. \url{https://www.semanticscholar.org/paper/46f9f7b8f88f72e12cbdb21e3311f995eb6e65c5}
\bibitem{ref3} Zhengbao Jiang, Frank F. Xu, Luyu Gao, Zhiqing Sun, Qian Liu, Jane Dwivedi-Yu, et al.. Active Retrieval Augmented Generation. \emph{Conference on Empirical Methods in Natural Language Processing}, 2023. \url{https://www.semanticscholar.org/paper/88884b8806262a4095036041e3567d450dba39f7}
\bibitem{ref4} Xin Cheng, Xun Wang, Xingxing Zhang, Tao Ge, Si-Qing Chen, Furu Wei, et al.. xRAG: Extreme Context Compression for Retrieval-augmented Generation with One Token. \emph{Neural Information Processing Systems}, 2024. \url{https://www.semanticscholar.org/paper/38fcc3667a907d6c94267c674aad114aae68441e}
\bibitem{ref5} ES Shahul, J. James, Luis Espinosa Anke, S. Schockaert. RAGAs: Automated Evaluation of Retrieval Augmented Generation. \emph{Conference of the European Chapter of the Association for Computational Linguistics}, 2023. \url{https://www.semanticscholar.org/paper/f5e9e5bbe22f0263be1f1ce88c66978a2b927772}
\bibitem{ref6} Yixuan Tang, Yi Yang. MultiHop-RAG: Benchmarking Retrieval-Augmented Generation for Multi-Hop Queries. \emph{arXiv.org}, 2024. \url{https://www.semanticscholar.org/paper/4e71624e90960cb003e311a0fe3b8be4c2863239}
\bibitem{ref7} Xiaoxin He, Yijun Tian, Yifei Sun, N. Chawla, T. Laurent, Yann LeCun, et al.. G-Retriever: Retrieval-Augmented Generation for Textual Graph Understanding and Question Answering. \emph{Neural Information Processing Systems}, 2024. \url{https://www.semanticscholar.org/paper/a41d4a3b005c8ec4f821e6ee96672d930ca9596c}
\bibitem{ref8} Zirui Guo, Lianghao Xia, Yanhua Yu, Tu Ao, Chao Huang. LightRAG: Simple and Fast Retrieval-Augmented Generation. \emph{Conference on Empirical Methods in Natural Language Processing}, 2024. \url{https://www.semanticscholar.org/paper/1ea143c34b9bc359780f79ba4d68dee68bcc1129}
\bibitem{ref9} Shijie Wang, Wenqi Fan, Yue Feng, Xinyu Ma, Shuaiqiang Wang, Dawei Yin. Knowledge Graph Retrieval-Augmented Generation for LLM-based Recommendation. \emph{Annual Meeting of the Association for Computational Linguistics}, 2025. \url{https://www.semanticscholar.org/paper/3ed128b6f0e08294fd9cb413eac96b4319481c67}
\bibitem{ref10} Siru Liu, Allison B. McCoy, Adam Wright. Improving large language model applications in biomedicine with retrieval-augmented generation: a systematic review, meta-analysis, and clinical development guidelines. \emph{J. Am. Medical Informatics Assoc.}, 2025. \url{https://www.semanticscholar.org/paper/83939671534dc3d374c9bc4e3e03b5ec2c7ba301}
\bibitem{ref11} L. M. Amugongo, Pietro Mascheroni, Steve Brooks, Stefan Doering, Jan Seidel. Retrieval augmented generation for large language models in healthcare: A systematic review. \emph{PLOS Digital Health}, 2025. \url{https://www.semanticscholar.org/paper/5879575701b9b65b5cc56c00d9eebbfa219e0428}
\bibitem{ref12} Guangzhi Xiong, Qiao Jin, Zhiyong Lu, Aidong Zhang. Benchmarking Retrieval-Augmented Generation for Medicine. \emph{Annual Meeting of the Association for Computational Linguistics}, 2024. \url{https://www.semanticscholar.org/paper/b798cf6af813638fab09a8af6ad0f3df6c241485}
\bibitem{ref13} Yuhe Ke, Liyuan Jin, Kabilan Elangovan, H. Abdullah, Nan Liu, Alex Tiong Heng Sia, et al.. Retrieval augmented generation for 10 large language models and its generalizability in assessing medical fitness. \emph{npj Digital Medicine}, 2025. \url{https://www.semanticscholar.org/paper/8c85026fe0d5c96fb275b81d483f8910db6ada03}
\bibitem{ref14} Xuejiao Zhao, Siyan Liu, Su-Yin Yang, C. Miao. MedRAG: Enhancing Retrieval-augmented Generation with Knowledge Graph-Elicited Reasoning for Healthcare Copilot. \emph{The Web Conference}, 2025. \url{https://www.semanticscholar.org/paper/da83852315c884c73dc527a4b7bc1209fbb037c3}
\bibitem{ref15} Wei Zou, Runpeng Geng, Binghui Wang, Jinyuan Jia. PoisonedRAG: Knowledge Corruption Attacks to Retrieval-Augmented Generation of Large Language Models. \emph{USENIX Security Symposium}, 2024. \url{https://www.semanticscholar.org/paper/f4e06256ab07727ff4e0465deea83fcf45012354}
\bibitem{ref16} A. Shafran, R. Schuster, Vitaly Shmatikov. Machine Against the RAG: Jamming Retrieval-Augmented Generation with Blocker Documents. \emph{USENIX Security Symposium}, 2024. \url{https://www.semanticscholar.org/paper/320b86d8a23da014d018db23a5f43b94368dd749}
\bibitem{ref17} Shenglai Zeng, Jiankun Zhang, Pengfei He, Yue Xing, Yiding Liu, Han Xu, et al.. The Good and The Bad: Exploring Privacy Issues in Retrieval-Augmented Generation (RAG). \emph{Annual Meeting of the Association for Computational Linguistics}, 2024. \url{https://www.semanticscholar.org/paper/ea89b058ce619ed16d4de633126b02a8179457c8}
\bibitem{ref18} Qinggang Zhang, Shengyuan Chen, Yuan-Qi Bei, Zheng Yuan, Huachi Zhou, Zijin Hong, et al.. A Survey of Graph Retrieval-Augmented Generation for Customized Large Language Models. \emph{arXiv.org}, 2025. \url{https://www.semanticscholar.org/paper/908d45b0d2b88ba72ee501c368eb618d29d61ce0}
\bibitem{ref19} Penghao Zhao, Hailin Zhang, Qinhan Yu, Zhengren Wang, Yunteng Geng, Fangcheng Fu, et al.. Retrieval-Augmented Generation for AI-Generated Content: A Survey. \emph{Data Science and Engineering}, 2024. \url{https://www.semanticscholar.org/paper/ab15463babf98fffc6f683fe2026de0725b5e1a9}
\bibitem{ref20} Aditi Singh, Abul Ehtesham, Saket Kumar, T. T. Khoei. Agentic Retrieval-Augmented Generation: A Survey on Agentic RAG. \emph{arXiv.org}, 2025. \url{https://www.semanticscholar.org/paper/ba7952e7c4fb891c36980ca19f94251257da6eb7}
\end{thebibliography}
\end{document}
